\chapter{Conclusione e Sviluppi Futuri}
\label{chap:conclusione}

Il lavoro svolto durante il tirocinio e descritto in questa tesi ha affrontato una delle limitazioni architetturali più significative del \textit{Project Emerge}: la rigida dipendenza dal protocollo \ac{MQTT} come unico mezzo di comunicazione tra i moduli del sistema. Questa dipendenza, seppur funzionale per gli obiettivi dimostrativi iniziali, rappresentava un ostacolo all'integrazione di robot e client eterogenei, limitando la flessibilità e l'estendibilità del sistema in contesti reali e dinamici.

Il contributo principale di questo lavoro è stato la progettazione e l'implementazione di un \textit{domain gateway} che funge da ponte tra il mondo interno del sistema, basato su \ac{MQTT}, e il mondo esterno, caratterizzato da protocolli eterogenei come \ac{HTTP} \ac{REST}, WebSocket e \ac{CoAP}. L'architettura del gateway si basa su tre pilastri fondamentali: un sistema di \textit{connessioni} modulari, ciascuna dedicata a un protocollo specifico; un meccanismo di \textit{bus interni} per lo scambio di messaggi tra le connessioni; e una \textit{cache} che garantisce la coerenza e la disponibilità immediata dei dati per i client che utilizzano modelli di interazione request/response. Questa struttura, combinata con l'uso della \textit{dependency injection}, ha permesso di ottenere un sistema estendibile, in cui l'aggiunta di un nuovo protocollo richiede esclusivamente l'implementazione di una nuova connessione, senza modificare il nucleo del gateway né i moduli preesistenti.

I requisiti funzionali e non funzionali sono stati pienamente soddisfatti: il gateway supporta la lettura e la scrittura delle risorse del sistema attraverso protocolli eterogenei, garantendo sia un modello di interazione a notifica continua (analogo al publish/subscribe di \ac{MQTT}) sia un modello request/response. Inoltre, l'adozione di un architettura asincrona basata su \texttt{asyncio} e l'uso di librerie specializzate come \texttt{FastAPI}, \texttt{aiomqtt} e \texttt{aiocoap} hanno assicurato prestazioni elevate anche in presenza di un numero significativo di connessioni simultanee. La suite di test automatici, realizzata con \texttt{pytest} e \texttt{pytest-asyncio}, ha permesso di validare il comportamento del gateway con una copertura dell'88\% del codice, garantendo affidabilità e manutenibilità.

Tuttavia, il lavoro presenta alcune limitazioni che aprono la strada a possibili evoluzioni future. In primo luogo, la copertura dei test non include la logica di riconnessione del modulo \ac{MQTT}, che richiederebbe un ambiente di test in grado di simulare un broker irraggiungibile. Inoltre, le prestazioni del gateway non sono state misurate quantitativamente in scenari di carico estremo, un aspetto che potrebbe essere approfondito in futuri lavori. Infine, l'attuale implementazione della cache è in-memory e non distribuita: in uno scenario con più istanze del gateway, sarebbe utile valutare l'integrazione con un sistema di caching distribuito come Redis.

Per quanto riguarda gli sviluppi futuri, diverse direzioni possono essere esplorate. Una delle più rilevanti è la sostituzione graduale di \ac{MQTT} come protocollo interno: l'architettura del gateway, grazie al disaccoppiamento tra il mondo interno ed esterno, permetterebbe di introdurre un nuovo protocollo interno senza modificare le connessioni esterne né i moduli esistenti di \textit{Emerge}. Un'altra possibilità è l'estensione del gateway con il supporto a ulteriori protocolli, come \ac{AMQP} o \ac{gRPC}, per rispondere a esigenze ancora più eterogenee. Infine, sarebbe interessante valutare l'introduzione di meccanismi di bilanciamento del carico e replicazione del gateway per garantire scalabilità orizzontale e alta disponibilità in ambienti di produzione.

In conclusione, questo lavoro ha dimostrato come un'attenta analisi dei vincoli architetturali e una progettazione modulare possano trasformare un sistema inizialmente limitato in una piattaforma flessibile e estendibile, pronta ad adattarsi a nuove esigenze e contesti applicativi. Il \textit{domain gateway} non solo risolve il problema immediato della dipendenza da \ac{MQTT}, ma fornisce anche una base solida per future evoluzioni del progetto \textit{Emerge} e, più in generale, per l'applicazione dell'\textit{Aggregate Computing} in scenari reali e complessi.